\section{The local reserve and marked one-body invitation}
\label{sec:ex-marked-field-local}

This unit closes two finite estimates used when a marked label is inserted
at the terminal scale of a selected history. The first is the smallness of
the forced reserve relative to a fixed section and its finite frame. The
second converts a paid local section into the exact coefficient and
remainder data called a \emph{Cream}. It keeps the actual one-body
insertion, its normalization, and its exponential support price visible.
The later recentering of this Cream through all earlier scales is a separate
step and is not claimed here.

\subsection{The forced reserve beats every fixed frame price}

Let $F_d$ be a finite frame with grades $(r_j,s_j)$, $j\in J_d$, and set
\[
 M_d=|J_d|,\qquad k_d=\max_{j\in J_d}s_j.
\]
For a budget threshold $\vartheta_N\in(0,1]$, define the finite seat
price
\[
 \operatorname{Seat}_{d,N}=
       \sum_{j\in J_d}\vartheta_N^{-s_j}.
\]
Assume the scale parameter $\beta_N\to\infty$ and
$\vartheta_N^{-1}\leq\beta_N$. Define
\[
 \operatorname{stain}(x)=
 \exp\!\left(-\exp\!\left(\frac{(\log\log x)^2}{5}\right)\right)
 \quad(x>1),\qquad
 \mathfrak f_N=\operatorname{stain}(\beta_N),
\]
and, for fixed $x,y\geq0$,
\[
 \operatorname{Storm}_{d,N}(x,y)=(x+2y)\operatorname{Seat}_{d,N}.
\]

\begin{lemma}[Reserve versus a fixed section]
\label{lem:ex-marked-sanded}
For every fixed $d$ and $x,y\geq0$,
\[
 \sqrt{\mathfrak f_N}\,\operatorname{Storm}_{d,N}(x,y)\leq1
\]
for all sufficiently large $N$.
For every fixed finite stencil $X$ on the cover lattice,
\[
 |\operatorname{regional}(X,0)|\,\mathfrak f_N\leq1
\]
for all sufficiently large $N$ as well.
\end{lemma}
\begin{proof}
Write $M=M_d$, $k=k_d$, and $A=x+2y$. Since each $s_j\leq k$ and
$\vartheta_N^{-1}\leq\beta_N$,
\[
 \operatorname{Seat}_{d,N}\leq M\beta_N^k.
\]
Choose $N$ large enough that $\beta_N\geq\max\{1,AM\}$. Then
$A M\beta_N^k\leq\beta_N^{k+1}$, so it suffices that
$\sqrt{\operatorname{stain}(\beta_N)}\,\beta_N^{k+1}\leq1$.
For completeness, put $v=\log\log x$. The logarithm of the left side is
\[
 (k+1)e^v-\tfrac12e^{v^2/5}.
\]
As $v\to\infty$, $v^2/5-v\to\infty$, hence
$e^{v^2/5}\geq2(k+1)e^v$ eventually. This proves the required
super-polynomial decay and the first assertion.

Let $m_X=|X|$. Regionalization is an image, so
$|\operatorname{regional}(X,0)|\leq m_X$. Eventually
$\beta_N\geq\max\{1,m_X\}$. The just-proved estimate with exponent one
implies $\sqrt{\mathfrak f_N}\,\beta_N\leq1$, while
$0<\mathfrak f_N\leq1$; therefore
\[
 |\operatorname{regional}(X,0)|\mathfrak f_N
 \leq m_X\sqrt{\mathfrak f_N}
 \leq \beta_N\sqrt{\mathfrak f_N}\leq1.
\]
This proves the finite local allowance estimate.
\end{proof}

The volume-independent geometric cover estimate is also elementary. Let
$B_N$ be the radius in the budget's support norm and assume $B_N\to\infty$.
For fixed $X$, every point of $X$ is eventually in the radius-$B_N$ cube.
Consequently $\operatorname{regional}(X,0)$ lies in the halo of the
origin. The cube-tiling cover of that halo uses at most $81$ support
blocks, so
\[
 \operatorname{coverNumber}(B_N,\operatorname{regional}(X,0))\leq81
 \label{eq:ex-marked-cover}
\]
eventually. In the source schedule $B_N\to\infty$ follows from
$\beta_N\to\infty$ and the radius estimate in the finite-volume budget.
The three local conditions used below are therefore obtained with the
same terminal threshold: the regional forced mass is at most one, its
support cover number is at most $81$, and the square-root reserve times
the fixed storms price is at most one.

\subsection{From a paid section to a Cream}

We record the exact one-body statement. Let $C$ be a selected finite-volume
family, $N$ its terminal level, $p$ the period, and $d$ a label degree.
Write $F=F_d$, $T$ for the fixed local section, $V$ for the purchased
section, $\theta_N\in(0,1]$ for the terminal budget threshold,
$w_N=\mathfrak f_N$ for the forced allowance, $u_N>0$ for the startup
allowance, and $\zeta_{d,N}\geq0$ for the mark budget. Let
$R_N=C.\mathrm{path.rep}(N,p,N)$ be the terminal representative,
$\alpha_N$ its support exponent, $a_N>0$ the physical lattice spacing, and
$B_N$ the support radius in the terminal budget point. Put
\begin{align*}
 X&=\operatorname{stencil}(V)\cup\operatorname{stencil}(T),\\
 D&=\operatorname{regional}(X,0),\qquad
 c_N=\text{column count},\\
 s_N&=\operatorname{Storm}_{d,N}\bigl(\operatorname{price}(V),
                                       \operatorname{price}(T)\bigr).
\end{align*}
where the two prices are the global value-norm constants of $V$ and $T$;
they are distinct from the derivative credit $b_N$. Assume
$\alpha_N\geq0$ and the following exact local inputs:

\begin{enumerate}
\item $C$ is a selected family which makes the action and
  $C.\mathrm{ref}(d)\leq N$. The section $V$ is sound, and its Credit holds
  at the terminal order and power with a nonnegative price $b_N$.
\item Every $z\in X$ satisfies
  $20(|z_j|+1)\leq\mathrm{side}(N,p)$ for each coordinate $j$.
  The local regional bounds are
  $0<w_N$, $|D|w_N\leq1$, and
  $\operatorname{coverNumber}(B_N,D)\leq81$.
\item The selected terminal representative is lawful and has a terminal
  Keeper. Also $u_N>0$, $\zeta_{d,N}\geq0$, and
  $\sqrt{w_N}s_N\leq1$.
\item At the terminal budget point, every placed basis column of $T$ is
  regular on its footprint and has derivative bound
  $c_N\theta_N^{s_j}$ for grade $s_j$; this is the selected-family price
  clause.
\end{enumerate}

For $v\in F$, define the normalized cylinder of the purchased section by
\[
 \mathcal O_{N,p}(v)(U)
   =(a_N^d)^{-1}V.\mathrm{places}(F.\mathrm{scaled}(\theta_N^{-1})v,0,U).
\]
For a basis vector $e_j$ of $F$, define
\begin{align*}
 Q_N&=2+e^{81\alpha_N}\bigl(b_N+c_N+2|X|\bigr),\\
 E_N&=4(1+c_N)Q_N
       +e^{81(\alpha_N+2)}s_N\frac{\zeta_{d,N}}{u_N}.
\end{align*}

\begin{proposition}[One-body marked invitation]
\label{prop:ex-marked-invitation}
Under these inputs, for each basis vector $e_j$ there is a coefficient
record $Z_j$ and a regular price $r_N=b_N+c_N$
such that, with
\begin{align*}
 t_N&=e^{81\alpha_N}r_N
       +2|D|\sqrt{w_N}\,e^{81\alpha_N}s_N,\\
 \overline A_N&=e^{81(\alpha_N+2)}s_N.
\end{align*}
The record fields are
\begin{align*}
 Z_j.x&=e_j,\\
 &Z_j.\mathrm{Format}(0,C.\mathrm{params},R_N),\\
 &Z_j.\mathrm{remBound}(0,R_N,\sqrt{w_N},t_N),\\
 F.\mathrm{coeffSize}(Z_j.x)+t_N&\leq Q_N.
\end{align*}
and $0<Q_N$. There are an attenuation $S_j$ and an actual reserved mass
\[
A_j=e^{(\alpha_N+2)\operatorname{coverNumber}(B_N,D)}s_N
\]
with
$\mathrm{Attainment}(R_N.p,0,S_j,A_j)$ and
$A_j\zeta_{d,N}\leq u_NE_N$, and
\[
 \mathrm{profile}(R_N.\mathrm{activities},Z_j.\mathrm{forced})
   =\mathrm{hook}(R_N.\mathrm{activities},R_N.p.\mathrm{threshold},S_j).
\]
For every configuration $U$, its density and continuity clauses are
\begin{align*}
 Z_j.\mathrm{density}(0,T,R_N,U)
   &=(a_N^d)R_N.\mathrm{eval}(U)b_j(U),\\
 b_j(U)&=(a_N^d)^{-1}
 V.\mathrm{places}\bigl(F.\mathrm{scaled}(\theta_N^{-1})e_j,0,U\bigr),\\
 &\hspace{25mm} b_j\text{ is continuous}.
\end{align*}
The reserved-mass bound is
\[
 \overline A_N\zeta_{d,N}\leq u_NE_N,
 \qquad 4(1+c_N)Q_N\leq E_N.
\]
Equivalently, these are the $\mathrm{Cream}$ fields used by the selected
field calibration.
\end{proposition}
\begin{proof}
Fix $j$. The coordinate vector
\[
 a=F.\mathrm{scaled}(\theta_N^{-1})e_j,
 \qquad z=F.\mathrm{coords}^{-1}(a).
\]
By the inverse-coordinate identity,
\[
          F.\mathrm{scaled}(\theta_N)\,F.\mathrm{coords}(z)=e_j.
\]
Set
$y=V.\mathrm{places}(a,0,\cdot)$. The pointwise section
$V.\mathrm{places}(a,0,\cdot)$ reads only the translated stencil of $V$;
by the stencil inclusion, this is contained in $D$. The margin assumption
implies that the placement is nonwrapping: each coordinate of every
stencil point has absolute value strictly smaller than one twentieth of
the torus side, so the translated stencil has a unique representative in
the centered box. Thus the source's read relation for $z$ holds on $D$.
The same statements apply to $T$.

The Credit estimate for $V$ applies to each basis column on its footprint.
For the rescaled axis $a=F.\mathrm{scaled}(\theta_N^{-1})e_j$, only column
$j$ is nonzero; its coefficient has modulus $\theta_N^{-s_j}$. The
terminal Credit gives every ordered derivative test admitted by the
terminal budget, on its healthy translated footprint, the bound
$b_N\theta_N^{s_j}$ for the unscaled column. Scaling cancels these factors,
so $y=V.\mathrm{places}(a,0,\cdot)$ has the same regularity and derivative
bound $b_N$. The selected-family price clause gives the identical
conclusion for $g=T.\mathrm{places}(a,0,\cdot)$ with bound $c_N$.
Regularity and derivative sup bounds are closed under subtraction, so
$r=y-g$ is regular on $D$ and every such derivative is bounded by
$r_N=b_N+c_N$. This is the bound used in the regular part of the
remainder. The global value-norm constants in $s_N$ are different: the
definition of $\operatorname{price}(V)$ and finite-coordinate expansion
give, for every coefficient vector $v$,
\[
 |V.\mathrm{places}(v,0,U)|
 \leq\operatorname{price}(V)F.\mathrm{coeffSize}(v),
\]
and likewise for $T$. The scaled-coefficient estimate gives the explicit
envelopes
\[
 \|V.\mathrm{places}(a,0,U)\|
       \leq\operatorname{price}(V)\operatorname{Seat}_{d,N},\qquad
 \|T.\mathrm{places}(a,0,U)\|
       \leq\operatorname{price}(T)\operatorname{Seat}_{d,N}.
\]
These are the uniform pointwise bounds used in the residual estimate. The local insertion
construction uses this read relation and the terminal Keeper to split the
target into its unit coordinate, the core of the local section, and a
residual which vanishes at the identity. With
$g=T.\mathrm{places}(a,0,\cdot)$ and residual $r=y-g$, define
\begin{align*}
 b_0&=r-\operatorname{mask}_{D}(r)
          +\xi-\operatorname{mask}_{A}(\xi),\\
 \xi&=T.\mathrm{places}\bigl(
       F.\mathrm{scaled}(\theta_N^{-1})F.\mathrm{core}(e_j),0,\cdot\bigr).
\end{align*}
where $A$ is the local-section footprint and each mask is the indicator of
the corresponding healthy set. The read identity, unit normalization, and
the projection condition in section soundness give the null-value and
gauge-vanishing conditions needed for this residual. On the healthy set for $D$,
$\operatorname{mask}_D(r)=r$ and, since $A\subseteq D$,
$\operatorname{mask}_A(\xi)=\xi$; hence $b_0=0$ there. The read and
projection identities give the required null-value and gauge-vanishing
conditions for $r$. The triangle inequality and the two section prices
then give the pointwise bound
\[
 |b_0(U)|\leq \operatorname{price}(V)\operatorname{Seat}_{d,N}
   +2\operatorname{price}(T)\operatorname{Seat}_{d,N}=s_N.
\]
This is the exact ``storms'' majorant; the two local-section terms are
charged separately.

The target entering the coefficient density is
\[
 b_j(U)=(a_N^d)^{-1}y(U),\qquad
 y(U)=V.\mathrm{places}(F.\mathrm{scaled}(\theta_N^{-1})e_j,0,U).
\]
Thus $b_j=\mathcal O_{N,p}(e_j)$.

For clarity, the record is built by the source's finite forcing
construction, not by an unspecified existence argument. Start with the
unpriced terminal record and its coefficient vector $e_j$. Its regular
part is the section remainder on $D$; the forced correction is supported
on $D$. To price the correction, let
\[
 \operatorname{mergePrice}(D,h)=e^{\alpha_N\operatorname{coverNumber}(B_N,D)}
 h\bigl(e^{|D|w_N}-1\bigr).
\]
When $|D|w_N\leq1$, the elementary inequality
$e^t-1\leq2t$ for $0\leq t\leq1$ yields
\[
 \operatorname{mergePrice}(D,h)
 \leq \sqrt{w_N}\,2|D|\sqrt{w_N}
       e^{\alpha_N\operatorname{coverNumber}(B_N,D)}h.
\]
Apply activity garnishing to the singleton family consisting of support
$D$ and correction $b_0$, using the canonical support, inventory, and
admissibility dictionary supplied by the terminal Keeper constructed in
Section~\ref{sec:terminal-marked-construction}. In the
construction, the temporary enlarged eligibility dictionary is transported
back to the original one by the canonical crossing map. The inputs are the
origin support, attachment, nonnegative bound $s_N$, and the merge-price
budget above. It returns a
forced family whose gas-mark sum is $b_0(U)$ times the activity gas density
and whose profile is the prescribed hook. Add this family to the regular
coefficient produced by the support split. Pointwise, the exponential entry
contributes the unit coordinate plus the masked residual, while the
gas-mark sum contributes $b_0$; their sum is the target. The regular
remainder's price is its ordinary mass plus the merge price. This
construction also gives an attainment with reserved mass
\[
A_j=e^{(\alpha_N+2)\operatorname{coverNumber}(B_N,D)}s_N
       \leq e^{81(\alpha_N+2)}s_N.
\]
The regular part of the remainder is bounded by
$e^{81\alpha_N}r_N$. The merge-price estimate, in the
$\sqrt{w_N}$-weighted remainder budget, contributes at most
$2|D|\sqrt{w_N}e^{81\alpha_N}s_N$. Their sum is $t_N$. Since
$r_N=b_N+c_N$, $|D|\leq|X|$, and
$\sqrt{w_N}s_N\leq1$,
\[
 t_N\leq e^{81\alpha_N}(b_N+c_N+2|X|),
 \qquad
 1+t_N\leq 2+e^{81\alpha_N}(b_N+c_N+2|X|)=Q_N.
\]
Here the extra $1$ is the coefficient size of $e_j$; this proves the
Cream charge clause. The actual reserved mass is
\[
A_j=e^{(\alpha_N+2)\operatorname{coverNumber}(B_N,D)}s_N,
\]
and the cover bound gives $A_j\leq\overline A_N$. Multiplying the latter envelope by
$\zeta_{d,N}/u_N$ gives the reserved-mass clause from the second summand
of $E_N$; the first summand gives $4(1+c_N)Q_N\leq E_N$. This proves all
the stated record fields.
\end{proof}

The proposition closes the one-body invitation consumed by
$\texttt{Parterre.inviting}$ once the exact terminal margins, Keeper,
Credit, and local regional bounds are in place. The remaining selected
field calibration step is \texttt{Figment.figment}: it chooses a single
recentered label for all earlier-scale root tests. Its source proof uses
the selected Music's ordered adjacent-scale transitions, Antilles'
reverse induction, the paid height/Caroling estimates, and the finite
cutoff remainder transfer. Those are not consequences of the local
reserve and invitation arguments above; their exact inputs and proof
status are tracked in the evidence ledger.
